% =============================================================================
% Project : THE_PHILOSOPHY_OF_VIABILITY
% Book    : Book I — Ontology of Viability
% Book ID : VS.B01
%
% File    : 4_ASSETS/diagrams/VS.B01.P01.C03.S06.SS10.P004.DG01.multiplicity_across_ontological_levels.tex
%
% Asset:
%   Diagram DG01 — Multiplicity Across Ontological Levels
%
% Purpose :
%   To distinguish possible subjects and levels at which multiplicity
%   may be examined, without presenting them as a necessary hierarchy
%   or inferring identity plurality from organizational plurality.
%
% =============================================================================

\begin{figure}[htbp]
\centering
\begin{tikzpicture}[
    node distance=1.15cm,
    every node/.style={align=center},
    levelnode/.style={
        draw,
        rounded corners,
        text width=0.69\textwidth,
        minimum height=0.85cm,
        inner sep=5pt
    },
    relation/.style={-{Stealth[length=2mm]}, thick, dashed}
]

\node[levelnode] (elements) {
    \textbf{Elements or determinations}\\
    \footnotesize
    Components, features, relations, manifestations, or processes
};

\node[levelnode, below=of elements] (subject) {
    \textbf{A specified subject}\\
    \footnotesize
    Existent, composition, aggregate, or configuration
};

\node[levelnode, below=of subject] (configuration) {
    \textbf{A broader configuration}\\
    \footnotesize
    Relations among subjects or participation in a wider arrangement
};

\node[
    draw,
    rounded corners,
    text width=0.69\textwidth,
    minimum height=1.0cm,
    inner sep=5pt,
    below=1.35cm of configuration
] (identity) {
    \textbf{Identity question}\\
    \footnotesize
    Unity, plurality, or unresolved identity relations
};

\draw[relation] (elements) -- node[right, font=\scriptsize, text width=2.2cm] {may belong to or characterize} (subject);
\draw[relation] (subject) -- node[right, font=\scriptsize, text width=2.2cm] {may participate in} (configuration);
\draw[relation] (configuration) -- node[right, font=\scriptsize, text width=2.2cm] {does not decide} (identity);

\node[
    text width=0.83\textwidth,
    align=center,
    font=\footnotesize,
    below=0.35cm of identity
] {
    The diagram distinguishes analytical questions. It does not assert
    that every configuration has an identity, that identity is located
    at a particular level, or that the levels form a necessary hierarchy.
};

\end{tikzpicture}
\caption{Multiplicity Across Ontological Levels}
\label{fig:multiplicity-across-ontological-levels}
\end{figure}

% =============================================================================
% END OF DIAGRAM DG01
% =============================================================================